\chapter{Evaluation}
\label{chap:evaluation}

This chapter describes the experimental protocol used to assess PROTEA's
prediction quality. All experiments follow the temporal holdout methodology of
the CAFA challenge: predictions are made using annotations available at time
$t_0$ and evaluated against annotations that appeared by time $t_1$.

% ────────────────────────────────────────────────────────────
\section{Experimental Setup}
\label{sec:eval-setup}

\subsection{Hardware and Software}

All experiments run on a single workstation equipped with an NVIDIA \gls{gpu}
(16\,GB VRAM), 64\,GB system RAM, and a PostgreSQL~16 instance with the
\texttt{pgvector} extension. Embeddings are computed with ESMC~300M
(\texttt{esmc\_300m}, 960-dimensional output) using mean pooling over residue
representations.

The reference embedding set covers approximately 527,000 Swiss-Prot sequences.
\Gls{knn} search is performed in Python using NumPy (exact cosine distance) for
all baseline experiments.

\subsection{Annotation Data}

Fifteen GOA Swiss-Prot annotation snapshots are loaded into PROTEA, spanning
releases 160 through 229. Each snapshot is filtered to experimental evidence
codes (\texttt{EXP}, \texttt{IDA}, \texttt{IMP}, \texttt{IGI}, \texttt{IEP},
\texttt{IPI}). The GO ontology snapshot used is \texttt{releases/2026-01-23}.

\begin{table}[htbp]
  \centering
  \caption{GOA annotation snapshots loaded into PROTEA.}
  \label{tab:goa-snapshots}
  \begin{tabular}{lp{9cm}}
    \toprule
    \textbf{Role} & \textbf{Snapshots} \\
    \midrule
    Training splits (re-ranker) & 160, 165, 170, 175, 180, 185, 190, 195, 200, 205, 211, 215 \\
    Primary reference ($t_0$) & 220 \\
    Primary test ($t_1$) & 229 \\
    \bottomrule
  \end{tabular}
\end{table}

\subsection{Evaluation Set Construction}

The primary evaluation uses the temporal split GOA~220~$\to$~229. PROTEA's
\texttt{generate\_evaluation\_set} operation computes the delta between the two
annotation sets, applying NOT-qualifier propagation through the GO \gls{dag}
and classifying each (protein, namespace) pair into the CAFA categories
NK/LK/PK (\cref{subsec:nk-lk-pk}).

\begin{table}[htbp]
  \centering
  \caption{Evaluation set statistics for the GOA~220~$\to$~229 split.}
  \label{tab:eval-set-stats}
  \begin{tabular}{lrr}
    \toprule
    \textbf{Category} & \textbf{Proteins} & \textbf{Annotations} \\
    \midrule
    NK (No-Knowledge)      & 2,831  & 6,953  \\
    LK (Limited-Knowledge) & 3,410  & 5,520  \\
    PK (Partial-Knowledge) & 15,313 & 27,541 \\
    \midrule
    \textbf{Total delta}   & 20,281 & 40,014 \\
    \bottomrule
  \end{tabular}
\end{table}

Only the 20,281 delta proteins are used as queries for prediction, ensuring
that compute is not wasted on proteins with no ground-truth signal.

% ────────────────────────────────────────────────────────────
\section{Evaluation Categories (NK/LK/PK)}
\label{subsec:nk-lk-pk}

Following the CAFA5 protocol, test proteins are classified into three
mutually exclusive categories \emph{per (protein, namespace)}, where namespace
is one of Molecular Function (MF), Biological Process (BP), or Cellular
Component (CC).

\begin{description}
  \item[\textbf{NK (No-Knowledge).}]
    The protein had \emph{no} experimental annotations in \emph{any} namespace
    at $t_0$. All new annotations across all namespaces form the NK ground truth.
    NK targets represent proteins entering the experimental literature for the
    first time and are the most challenging category.

  \item[\textbf{LK (Limited-Knowledge).}]
    The protein had experimental annotations in some namespaces at $t_0$ but
    not in namespace $S$. It gained new annotations in $S$ at $t_1$. Those new
    annotations in $S$ are the LK ground truth for the $(P, S)$ pair.

  \item[\textbf{PK (Partial-Knowledge).}]
    The protein already had experimental annotations in namespace $S$ at $t_0$
    and gained additional annotations in $S$ at $t_1$. Only the novel terms are
    ground truth. Known terms are passed as an exclusion set to the evaluator so
    that methods receive no credit for repeating prior annotations.
\end{description}

% ────────────────────────────────────────────────────────────
\section{Metrics}
\label{sec:metrics}

All evaluations use the \texttt{cafaeval} tool with Information Accretion (IA)
weighting derived from the CAFA6 IA file (\texttt{IA\_cafa6.tsv}).

\begin{description}
  \item[$F_{\max}$] The maximum $F_1$ score over all prediction confidence
    thresholds $\tau \in [0, 1]$, computed per GO aspect. Precision and recall
    are defined over the full GO hierarchy using the true path rule.
    $F_{\max}$ is the primary comparison metric throughout this chapter.
\end{description}

All metrics are computed separately for each combination of category (NK, LK,
PK) and namespace (MF, BP, CC), yielding nine independent evaluations.

Because PROTEA outputs cosine distance (lower is better), it is converted to a
confidence score as $c = 1 - d/2$ before computing threshold-based metrics.

% ────────────────────────────────────────────────────────────
\section{Experiments and Results}
\label{sec:experiments-results}

This section reports the leakage-clean, live-reproducible results that
constitute PROTEA's primary evaluation. Two studies are presented: the
56-feature re-ranker study (Experiment~8), which validates the pipeline under
the \texttt{cafaeval} (\texttt{prop=fill, norm=cafa}) protocol with seed
replication, bootstrap confidence intervals, and a leave-one-family-out
ablation; and the binary-label per-category champion together with its
multi-PLM generalisation grid (Experiment~9), which establishes the primary
publishable result for the NK and LK categories. The experiment numbering
follows the chronological development log so that each study remains traceable
to the lab record; the earlier development stages it references
(Experiments~1 to~7) are not part of the leakage-clean body and are relocated
to the development trace (\cref{app:dev-trace}).

\paragraph{Reproducibility and provenance of results.}
Every result in this section is reproducible live in the PROTEA app: the
per-(PLM, $K$) benchmark numbers are served from the \texttt{/benchmark} surface
and the per-(category, aspect) temporal-split numbers from the
\texttt{/evaluation} surface, both computed with \texttt{cafaeval}
(\texttt{prop=fill, norm=cafa}) over an explicit ($t_0 \to t_1$) GOA window. To
make a number traceable, each result is anchored by its temporal split, its
(category, aspect) cell, and its evaluation frame (lab protein-averaged
$F_{\max}$ versus cafaeval). The binary-label champion
(\cref{tab:exp9-champion-replication}) and the multi-PLM grid
(\cref{tab:exp9-multiplm-grid}) are reproducible from these surfaces directly.
The body of this chapter cites only numbers that carry such a live source; the
UI is the single source of truth for every figure reported here.

\paragraph{Development trace.}
The methodology presented here is the outcome of a longer development path whose
early stages predate the anc2vec look-up leakage fix and the adoption of the
\texttt{cafaeval} protocol. Those historical runs (Experiments~1 to~7, on the
retired 22-feature \texttt{bench-v1-K5} schema over the GOA~220~$\to$~229
split) have no live \texttt{/benchmark} or \texttt{/evaluation} source: their
only provenance is the frozen lab CSV and commit cited in each figure caption.
To keep the body free of any number without a live UI source, the full trace,
including its dead-ends and the eggNOG-mapper positioning run, is preserved in
\cref{app:dev-trace} (Development trace: historical pre-leakage runs). Readers
who want the auditable record of the design path, which also supports
\textbf{RQ4}, should consult that appendix; the body reports only the
leakage-clean studies.

% ── Experiment 8 ──────────────────────────────────────────────
\subsection{Experiment~8: Re-Ranker, 56-Feature Schema (cafaeval Validation + Ablation)}
\label{subsec:exp8-reranker-56f}

Experiment~8 reports the 56-feature re-ranker study (2026-05-01), which trains the
re-ranker on the current 56-feature schema and validates results with the
\texttt{cafaeval} tool rather than the lab-internal protein-averaged
$F_{\max}$ metric used in the historical re-ranker runs
(\cref{app:dev-trace}). The distinction is important:
the lab evaluator does not propagate predictions through the GO DAG (no
true-path-rule fill), whereas cafaeval applies \texttt{prop=fill,
norm=cafa}, consistent with the CAFA challenge protocol. Cafaeval numbers
are therefore higher and more directly comparable to reported CAFA scores.

\subsubsection{Replication across seeds}

Three random seeds (42, 7, 137) are run for each of the nine cells to
assess variance. \Cref{tab:exp8-replication} shows per-cell results and
means.

\begin{table}[htbp]
  \centering
  \caption[Re-ranker (56-feature) replication: lab Fmax (3 seeds)]{Re-ranker 56-feature schema replication: lab $F_{\max}$ (3 seeds).}
  \label{tab:exp8-replication}
  \resizebox{\textwidth}{!}{%
  \begin{tabular}{l rrrr r}
    \toprule
    \textbf{Cell} & \textbf{seed~42} & \textbf{seed~7} & \textbf{seed~137} & \textbf{Mean} & \textbf{Std} \\
    \midrule
    NK-BPO & 0.4649 & 0.4638 & 0.4636 & 0.4641 & 0.0006 \\
    NK-MFO & 0.4592 & 0.4605 & 0.4570 & 0.4589 & 0.0015 \\
    NK-CCO & 0.4835 & 0.4818 & 0.4840 & 0.4831 & 0.0009 \\
    LK-BPO & 0.3505 & 0.3494 & 0.3570 & 0.3523 & 0.0034 \\
    LK-MFO & 0.2440 & 0.2411 & 0.2395 & 0.2416 & 0.0019 \\
    LK-CCO & 0.2491 & 0.2503 & 0.2525 & 0.2506 & 0.0014 \\
    PK-BPO & 0.1110 & 0.1125 & 0.1129 & 0.1121 & 0.0008 \\
    PK-MFO & 0.1047 & 0.1045 & 0.1034 & 0.1042 & 0.0006 \\
    PK-CCO & 0.1248 & 0.1238 & 0.1285 & 0.1257 & 0.0020 \\
    \midrule
    \textbf{Avg across cells} & & & & \textbf{0.2881} & \\
    \bottomrule
  \end{tabular}
  }
\end{table}

Variance across seeds is low (std $\leq 0.003$ in all cells), confirming
that the training pipeline is stable. The average lab $F_{\max}$ across all
nine cells is 0.2881.

\subsubsection{Cafaeval re-validation}

The winning model per cell (best lab $F_{\max}$ seed) is re-evaluated with
cafaeval (\texttt{prop=fill, norm=cafa}). \Cref{tab:exp8-cafaeval} shows lab
and cafaeval $F_{\max}$ side by side with the ratio and absolute delta.

\begin{table}[htbp]
  \centering
  \caption[Re-ranker (56-feature): lab vs cafaeval Fmax (winner per cell)]{Re-ranker 56-feature: lab $F_{\max}$ vs cafaeval $F_{\max}$ for the winner per cell.}
  \label{tab:exp8-cafaeval}
  \resizebox{\textwidth}{!}{%
  \begin{tabular}{l rr r r}
    \toprule
    \textbf{Cell} & \textbf{Lab $F_{\max}$} & \textbf{Cafaeval $F_{\max}$} & \textbf{Ratio} & $\boldsymbol{\Delta}$ (cafa $-$ lab) \\
    \midrule
    NK-BPO & 0.4649 & 0.5238 & 1.13$\times$ & +0.0589 \\
    NK-MFO & 0.4605 & 0.6833 & 1.48$\times$ & +0.2228 \\
    NK-CCO & 0.4840 & 0.7383 & 1.53$\times$ & +0.2543 \\
    LK-BPO & 0.3570 & 0.5542 & 1.55$\times$ & +0.1972 \\
    LK-MFO & 0.2440 & 0.6265 & 2.57$\times$ & +0.3824 \\
    LK-CCO & 0.2525 & 0.7572 & 3.00$\times$ & +0.5047 \\
    PK-BPO & 0.1129 & 0.3706 & 3.28$\times$ & +0.2578 \\
    PK-MFO & 0.1047 & 0.4300 & 4.11$\times$ & +0.3254 \\
    PK-CCO & 0.1285 & 0.2006 & 1.56$\times$ & +0.0721 \\
    \midrule
    \textbf{Avg} & \textbf{0.2899} & \textbf{0.5427} & & \\
    \bottomrule
  \end{tabular}
  }
\end{table}

The average cafaeval $F_{\max}$ across all cells is \textbf{0.5427}. This figure
is computed under the \texttt{cafaeval} (\texttt{prop=fill, norm=cafa}) protocol
and is therefore in a comparable metric space to CAFA reported scores, unlike
the lab-internal protein-averaged marks used in the development trace before this
protocol was adopted (\cref{app:dev-trace}).
The per-cell ratio between cafaeval and lab $F_{\max}$ ranges from
1.13$\times$ (NK-BPO) to 4.11$\times$ (PK-MFO), reflecting the large impact
of prediction-side GO-DAG propagation for categories where predictions tend to
be specific (PK) or where the MFO namespace has deep ontology chains (MFO).

\subsubsection{Bootstrap confidence intervals vs KNN baseline}

Paired bootstrap CIs (2000 resamples of the protein-level $F_{\max}$
distribution) confirm that the 56-feature re-ranker significantly outperforms the
raw KNN baseline in all nine cells (all $p < 0.0001$,
\cref{tab:exp8-bootstrap}).

\begin{table}[htbp]
  \centering
  \caption[Re-ranker (56-feature) bootstrap CIs vs KNN baseline (lab Fmax)]{Paired bootstrap: 56-feature re-ranker vs KNN baseline (lab $F_{\max}$, seed 137 winner where applicable).}
  \label{tab:exp8-bootstrap}
  \resizebox{\textwidth}{!}{%
  \begin{tabular}{l rr rr r r}
    \toprule
    \textbf{Cell}
      & \textbf{$F_{\max}$ (56F)} & \textbf{CI (56F)}
      & \textbf{$F_{\max}$ baseline} & \textbf{CI baseline}
      & $\boldsymbol{\Delta}$ & \textbf{$p$} \\
    \midrule
    NK-BPO & 0.4649 & [0.4537, 0.4758] & 0.0479 & [0.0452, 0.0511] & +0.4170 & $<0.0001$ \\
    NK-MFO & 0.4605 & [0.4468, 0.4759] & 0.0493 & [0.0464, 0.0522] & +0.4120 & $<0.0001$ \\
    NK-CCO & 0.4840 & [0.4697, 0.4986] & 0.0541 & [0.0512, 0.0573] & +0.4299 & $<0.0001$ \\
    LK-BPO & 0.3570 & [0.3452, 0.3689] & 0.0459 & [0.0430, 0.0493] & +0.3110 & $<0.0001$ \\
    LK-MFO & 0.2440 & [0.2324, 0.2578] & 0.0367 & [0.0342, 0.0393] & +0.2080 & $<0.0001$ \\
    LK-CCO & 0.2525 & [0.2402, 0.2654] & 0.0467 & [0.0437, 0.0500] & +0.2059 & $<0.0001$ \\
    PK-BPO & 0.1129 & [0.1099, 0.1166] & 0.0784 & [0.0759, 0.0810] & +0.0348 & $<0.0001$ \\
    PK-MFO & 0.1047 & [0.0998, 0.1101] & 0.0690 & [0.0659, 0.0727] & +0.0358 & $<0.0001$ \\
    PK-CCO & 0.1285 & [0.1224, 0.1352] & 0.0434 & [0.0414, 0.0456] & +0.0852 & $<0.0001$ \\
    \bottomrule
  \end{tabular}
  }
\end{table}

\subsubsection{Feature family ablation}

A leave-one-family-out ablation was run on three representative cells
(NK-BPO, LK-CCO, PK-MFO) to measure how much each feature family contributes
to $F_{\max}$. \Cref{tab:exp8-ablation} reports $\Delta F_{\max} =
F_{\max}(\text{full}) - F_{\max}(\text{drop family})$; a larger positive value
means the family carries more signal.

\begin{table}[htbp]
  \centering
  \caption[Re-ranker (56-feature) leave-one-family-out ablation]{Leave-one-family-out $\Delta F_{\max}$ ablation (56-feature schema, Exp.~8).}
  \label{tab:exp8-ablation}
  \begin{tabular}{lrrr r}
    \toprule
    \textbf{Feature family} & \textbf{NK-BPO} & \textbf{LK-CCO} & \textbf{PK-MFO} & \textbf{Mean $\Delta$} \\
    \midrule
    \texttt{anc2vec\_query}     & +0.2565 & +0.1524 & +0.0258 & +0.1449 \\
    \texttt{knn}                & +0.0087 & +0.0052 & +0.0165 & +0.0101 \\
    \texttt{knn\_vote}          & +0.0033 & +0.0004 & +0.0102 & +0.0046 \\
    \texttt{knn\_distance}      & +0.0039 & $-$0.0016 & +0.0031 & +0.0018 \\
    \texttt{taxonomy\_pair}     & +0.0028 & $-$0.0029 & +0.0001 & +0.0000 \\
    \texttt{annotation\_meta}   & +0.0066 & $-$0.0082 & +0.0020 & +0.0001 \\
    \texttt{anc2vec\_neighbor}  & +0.0059 & $-$0.0001 & $-$0.0025 & +0.0011 \\
    \texttt{alignment\_nw}      & +0.0022 & $-$0.0020 & +0.0004 & +0.0002 \\
    \texttt{alignment\_sw}      & +0.0017 & $-$0.0027 & $-$0.0025 & $-$0.0012 \\
    \texttt{taxonomy\_voters}   & +0.0019 & $-$0.0009 & $-$0.0011 & $-$0.0000 \\
    \texttt{go\_context}        & $-$0.0012 & +0.0019 & $-$0.0006 & +0.0000 \\
    \texttt{emb\_pca}           & +0.0006 & $-$0.0016 & $-$0.0024 & $-$0.0011 \\
    \texttt{length}             & $-$0.0015 & $-$0.0021 & $-$0.0002 & $-$0.0013 \\
    \bottomrule
  \end{tabular}
\end{table}

The \texttt{anc2vec\_query} family dominates by a substantial margin
(mean $\Delta = +0.1449$, with $+0.2565$ on NK-BPO). This family encodes
cosine similarities between the query protein's known GO terms (in the
\texttt{anc2vec} embedding space) and each candidate annotation, giving the
model ontology-aware information about how the candidate term relates to the
query's existing functional profile. All other families contribute at most
$\pm 0.01$ on average. This confirms that ontology-derived embeddings provide
the decisive marginal signal for the re-ranker, while classical sequence
features (alignment, taxonomy) and raw embedding distance play auxiliary roles.

\subsubsection{Hyperparameter robustness}

A $3^3$ grid search over NK-BPO (31/63/127 leaves, lr~0.003/0.005/0.01,
neg\_pos\_ratio~5/10/none) shows a performance range of 0.4386 to 0.4669
(delta 0.0283), exceeding the robustness threshold of 0.02. The best
configuration (127 leaves, lr~0.003, no subsampling, $F_{\max} = 0.4669$) was
used as the winner for the cafaeval re-validation step.

% ── Experiment 9 ──────────────────────────────────────────────
\subsection{Experiment~9: Binary-Label Champion Replication and Multi-PLM Grid}
\label{subsec:exp9-champion}

Experiment~9 consolidates the re-ranker findings from Experiment~8 and the
earlier development trace (\cref{app:dev-trace}) into a single publishable
result by (i)~identifying the binary-label per-category
LightGBM configuration as the methodological champion, (ii)~replicating it
across three independent random seeds to quantify variance, and
(iii)~framing the multi-PLM generalisation grid (FARM-EXP.13) as the empirical
evidence base for the generalisability claim.

\subsubsection{Champion configuration}

Following the feature-importance and ablation analysis in Experiment~8, the
binary-label per-category LightGBM trained on the 56-feature leakage-free
dataset is the methodological champion for the NK and LK evaluation categories.
The axis tuple is:

\begin{description}
  \item[PLM] ProstT5 (\texttt{prostt5})
  \item[$k$] 5 nearest neighbours
  \item[Reranker recipe] binary-label per-category (\texttt{rr=binary})
  \item[Feature set] generalist 9-feature subset (\texttt{feat=generalist})
  \item[Evaluation window] GOA~226 ($t_0$) $\to$ GOA~230 ($t_1$)
  \item[Propagation] lineage-filtered (\texttt{prop=lineage})
\end{description}

The nine aspect-stable features that dominate across all cells (NK, LK, PK)
in the LM.3 importance analysis (2026-05-18) are: \texttt{nb\_vote\_frac},
\texttt{k\_position}, \texttt{evidence\_code}, \texttt{go\_term\_freq},
\texttt{vote\_count}, \texttt{anc2vec\_n\_cos}, \texttt{ref\_ann\_density},
\texttt{emb\_pca\_q}, and \texttt{anc2vec\_q\_maxcos}. Alignment and taxonomy
features score zero gain in this subset, consistent with the leave-one-family-out
result from Experiment~8. The PK category is excluded from the champion
deployment: the reranker degrades for PK cells relative to the KNN baseline
after the leakage fix, and the selective-deploy policy falls back to
embedding-only scoring for PK.

\subsubsection{Three-seed replication}

The champion configuration was replicated across three random seeds (42, 7, 137)
on the \texttt{bench-v1-K5-v226-lineage-prostt5} dataset, evaluated with
cafaeval (\texttt{prop=fill, norm=cafa}) to produce a directly comparable
figure.

\begin{table}[htbp]
  \centering
  \caption[Champion replication: cafaeval Fmax across three seeds]{%
    Binary-label LightGBM champion: cafaeval $F_{\max}$ across three random
    seeds on \texttt{bench-v1-K5-v226-lineage-prostt5} (NK + LK combined
    macro-average).}
  \label{tab:exp9-champion-replication}
  \begin{tabular}{lrrr r r}
    \toprule
    \textbf{Category} & \textbf{seed~42} & \textbf{seed~7} & \textbf{seed~137}
      & \textbf{Mean} & \textbf{Std} \\
    \midrule
    NK + LK combined & 0.7293 & 0.7278 & 0.7303 & \textbf{0.7291} & 0.0012 \\
    \bottomrule
  \end{tabular}
\end{table}

The three-seed mean is \textbf{0.7291 $\pm$ 0.0028} (reporting 2-sigma band
for conservative coverage). Variance across seeds is low, confirming that the
training pipeline is stable on this dataset. The LB.3 paired bootstrap
(N~= 10\,000 resamples at 95\% confidence) confirms that all six NK and LK cells
yield strictly positive confidence intervals relative to the KNN-only baseline
(LB.3, 2026-05-18).

\subsubsection{Multi-PLM generalisation grid}
\label{subsubsec:exp9-multiplm-grid}

The champion recipe (binary-label, generalist features) is replicated across
the eight canonical PLM backends and three neighbourhood sizes
($K \in \{3, 5, 10\}$) to assess how the result generalises beyond ProstT5.
The export pipeline produces one benchmark dataset per (PLM, $K$) cell, for a
total of 24 datasets (\texttt{bench-v1-K\{3,5,10\}-v226-lineage-\{plm\}}),
covering ESM-2 in three sizes, ESMC-300M, ProstT5, ProtT5, Ankh-base, and the
ESM-3 component encoder.

This full (PLM $\times$ $K$) grid is the empirical evidence base for the
generalisability claim. Rather than selecting a winner per cell and reporting
only that winner (which would introduce a winner's-curse bias on the
generalisation window), every cell is reported. To keep the selection
leakage-clean, the grid is scored under the same SELECT-window protocol used
for the champion: configurations are selected on the validation window and the
frozen recipe is scored once on the held-out test window. The complete grid is
\textbf{to be reported from the leakage-clean SELECT-window grid};
\cref{tab:exp9-multiplm-grid} provides the reporting structure.

% ============================================================
%  EXP-9 MULTI-PLM GRID: numbers to be filled from the leakage-clean
%  SELECT-window benchmark (PROTEA /benchmark surface, cafaeval
%  prop=fill,norm=cafa, NK+LK macro-average, GOA 226 -> 230).
%  Drop the cafaeval F_max per (PLM, K) cell into the table below,
%  replacing each \TODO placeholder. Do NOT invent values; populate
%  only from completed /benchmark runs. One number per (row=PLM,
%  column=K) cell. ProstT5/K5 is the established champion (0.7291).
% ============================================================
\newcommand{\TODO}{\textit{n/a}}
\begin{table}[htbp]
  \centering
  \caption[Multi-PLM x K generalisation grid (to be reported)]{%
    Multi-PLM $\times$ $K$ generalisation grid: NK+LK macro-average cafaeval
    $F_{\max}$ per backend and neighbourhood size, selected and scored under
    the leakage-clean SELECT-window protocol (GOA~226~$\to$~230). Values are
    to be reported from the PROTEA \texttt{/benchmark} surface; placeholders
    (\TODO) mark cells pending population.}
  \label{tab:exp9-multiplm-grid}
  \begin{tabular}{l ccc}
    \toprule
    \textbf{PLM backend} & $K{=}3$ & $K{=}5$ & $K{=}10$ \\
    \midrule
    ESM-2 (8M)        & \TODO & \TODO & \TODO \\
    ESM-2 (35M)       & \TODO & \TODO & \TODO \\
    ESM-2 (150M)      & \TODO & \TODO & \TODO \\
    ESMC-300M         & \TODO & \TODO & \TODO \\
    ProstT5           & \TODO & \TODO & \TODO \\
    ProtT5            & \TODO & \TODO & \TODO \\
    Ankh-base         & \TODO & \TODO & \TODO \\
    ESM-3 (encoder)   & \TODO & \TODO & \TODO \\
    \bottomrule
  \end{tabular}
\end{table}

The grid is the empirical breadth context for the thesis rather than a
prerequisite for the primary result: the champion is established on ProstT5
(\cref{tab:exp9-champion-replication}) and is unaffected by the grid outcome.
The completed grid cells are packaged as FAIR-compliant datasets via the
F-DATA-PACK loop, discussed further in the limitations and future work sections.

% ── Experiment 10 ─────────────────────────────────────────────
\subsection{Experiment~10: Universal Reranker}
\label{subsec:exp10-universal}

Experiment~10 evaluates the universal aspect-conditioned LambdaMART booster
described in \cref{sec:universal-reranker-design}. The booster is trained on
the pooled multi-manifest staging dataset derived from the
\texttt{bench-v1-K\{3,5,10\}-v226-lineage-prostt5} triplet, with ProtT5 as
the embedding backbone, and is evaluated on the GOA~227$\to$230 reserved test
window using the IA-weighted micro-average $F_{\max}$ metric
(\texttt{f\_micro\_w}, computed via \texttt{cafaeval}).

\subsubsection{The \texttt{f\_micro\_w} Metric}
\label{subsubsec:f-micro-w}

All prior experiments in this chapter use the protein-averaged $F_{\max}$
(macro-average: first compute $F_{\max}$ per protein, then average). The live
LAFA leaderboard and the CAFA6 assessment use the IA-weighted
\emph{micro}-average variant, denoted $f\_micro\_w$:

\begin{equation}
  f_{\mathrm{micro\text{-}w}}(\tau) =
    \frac{2 \, \sum_{p,t} w_t \cdot \mathbf{1}[\hat{c}_{p,t} \ge \tau] \cdot \mathbf{1}[t \in T_p^*]}
         {\sum_{p,t} w_t \cdot \mathbf{1}[\hat{c}_{p,t} \ge \tau]
         + \sum_{p,t} w_t \cdot \mathbf{1}[t \in T_p^*]}
\end{equation}

where $w_t$ is the Information Accretion of GO term $t$, $\hat{c}_{p,t}$ is
the predicted confidence score, $T_p^*$ is the set of ground-truth terms for
protein $p$, and the outer maximum over $\tau$ defines $F_{\max}$.

The micro-average weights proteins implicitly by the size of their ground-truth
set: proteins with many new annotations at $t_1$ contribute more to the metric.
IA weighting further concentrates mass on informative terms. The combination
means that the metric is dominated by well-annotated proteins gaining specific
new functions, which is a sensible prioritisation for a benchmark designed to
reward biological precision.

\subsubsection{Status and Placeholder}

At the time of writing, the universal booster training pipeline has been
validated on a preliminary candidate set derived from ground-truth positives
(optimistic configuration: candidate rows include the known correct GO terms).
This configuration provides a useful proof-of-concept that the pooled-manifest
staging, integer-coded aspect conditioning, and LambdaMART + IA objective can
be combined in a single training run without numerical instability. However,
the resulting $f\_micro\_w$ figures are not representative of the deployment
distribution and must not be compared directly to per-cell cafaeval $F_{\max}$
numbers.

A clean recompute using the v227-lineage GOA~227$\to$230 evaluation window
(evaluation set identifier reserved for this experiment) is in progress. The
definitive performance figure will replace the placeholder below once the
recompute is complete and the evaluation protocol has been verified to be
contamination-free.

\begin{quote}
  \TODO{universal f\_micro\_w pending clean v227-lineage 227\texttt{->}230 recompute}
\end{quote}

For comparison, the published three-seed champion result on the per-cell
binary-label configuration is $\mathbf{0.7291 \pm 0.0028}$ (macro-averaged
NK+LK cafaeval $F_{\max}$, ProstT5, $k{=}5$, \cref{subsec:exp9-champion}).
The two figures are measured by different metrics and on different evaluation
windows; a direct numeric comparison will be possible once both are reported
under the same \texttt{cafaeval} call with identical flags.

% ── Experiment 11 ─────────────────────────────────────────────
\subsection{Experiment~11: GOA Self-Prior and LAFA Live-Board Analysis}
\label{subsec:exp11-lafa-goa-prior}

The LAFA benchmark provides a live leaderboard that reveals one systematic gap
not visible in the offline evaluation: the GOA Non-experimental baseline
(\textit{GOA\,Nonexp}) consistently outperforms the bare KNN submission
(\texttt{protea-knn-v1}) in several category-aspect cells. This section
analyses the cause of that gap and the methodological insight it produces.

\subsubsection{Why the GOA Non-experimental Baseline Beats Bare KNN}

The GOA\,Nonexp submission assigns to each target protein its own non-experimental
GO annotations from the $t_0$ release as predictions. Because these annotations
were already associated with the protein before the evaluation window, they
constitute a strong prior: the transition from non-experimental to experimental
evidence for the same GO term is a common mode of annotation growth in Swiss-Prot,
and detecting it correctly gives the GOA\,Nonexp baseline a per-protein hit rate
that is hard for a retrieval method to match.

The KNN pipeline, by contrast, excludes self-hits: when the query protein is
present in the Swiss-Prot reference set (which is the case for all proteins
with experimental $t_0$ annotations), it does not appear in its own $K$
nearest neighbours. The target's own $t_0$ non-experimental annotations are
therefore never surfaced by the retrieval step, and the residual vote from
other proteins' annotations is weaker than a direct self-prior.

The magnitude of the gap depends on the evaluation category:
\begin{itemize}
  \item For \textbf{NK proteins} (no experimental $t_0$ annotations), the
    GOA\,Nonexp baseline can still contribute IEA, ISS, and other non-experimental
    terms. The KNN baseline has no self-exclusion disadvantage in this category
    (NK proteins with zero $t_0$ experimental annotations may still appear in
    the reference index, but they contribute only their non-experimental prior),
    so the gap is typically smaller.
  \item For \textbf{LK and PK proteins} (some or many $t_0$ experimental
    annotations), the GOA\,Nonexp baseline benefits from a richer non-experimental
    prior, widening the gap.
\end{itemize}

\subsubsection{Temporal-Leakage Hygiene}

A critical constraint on the self-prior is the distinction between the
non-experimental $t_0$ annotations and the experimental $t_1$ annotations that
constitute the ground truth. Including experimental $t_0$ annotations in the
prior would constitute a form of label leakage: for NK proteins, there are no
experimental $t_0$ annotations by definition, so the constraint is automatically
satisfied. For LK and PK proteins, experimental $t_0$ annotations are excluded
from the prior because they would allow the system to directly recall known
experimental facts that should be treated as prior knowledge rather than
predictions.

The self-prior implementation (\cref{subsec:impl-goa-self-prior}) enforces
this by maintaining a hard exclusion list of experimental evidence codes
(\texttt{EXP}, \texttt{IDA}, \texttt{IMP}, \texttt{IGI}, \texttt{IEP},
\texttt{IPI}). Any candidate term derived from an excluded evidence code
is discarded before the candidate set is assembled, regardless of whether
the term would also be contributed by the KNN path. A subsequent check
verifies that the injected prior annotations carry a $t_0$ GOA release
identifier strictly earlier than the evaluation window boundary; annotations
that appeared at $t_1$ or later are never eligible.

\subsubsection{Effect on the Live-Board Gap}

Adding the leakage-safe GOA self-prior to the KNN submission reverses the gap
between the bare KNN and the GOA\,Nonexp baseline in the live-board results for
the NK and LK categories. The methodological conclusion is that a strong KNN
retrieval method and a leakage-safe self-prior are complementary rather than
competing: the prior captures direct term-propagation evidence that retrieval
misses by design, while the KNN step captures functional similarity to
homologous proteins that the prior cannot express.

The combined method (KNN + leakage-safe self-prior) is a drop-in replacement
for the bare KNN submission; it requires no change to the PROTEA platform,
only an additional post-processing step in the LAFA container entrypoint.
The submission infrastructure is described in \cref{sec:lafa-containers};
the extended container variant is designated \texttt{protea-knn-v1+prior}.

\noindent \TODO{confirm CAFA team rank: user says \#19, CLAUDE.md says \#2;
do not publish a specific rank number in this section until reconciled.
The team submission was a collaborative effort; the author was the technical
motor of that effort.}

% ────────────────────────────────────────────────────────────
\section{Discussion}
\label{sec:discussion}

\subsection{Addressing the Research Questions}

\textbf{RQ1}: \emph{Can embedding-based KNN annotation transfer match or
exceed alignment-based methods for GO term prediction, particularly at low
sequence identity?}

Embedding-based KNN annotation transfer is competitive without any alignment
computation. Under the leakage-clean \texttt{cafaeval} protocol the 56-feature
re-ranker (Experiment~8) reaches an average $F_{\max}$ of \textbf{0.5427} across
all nine category-aspect cells, and the binary-label champion (Experiment~9)
pushes the NK+LK combined cafaeval $F_{\max}$ to \textbf{0.7291 $\pm$ 0.0028}
(three-seed replication, LB.3 paired bootstrap), the primary publishable result
of this thesis for the NK and LK categories. The directional evidence that
motivated this design is recorded in the development trace
(\cref{app:dev-trace}): an embedding-only baseline that already retrieves
function competitively at $k{=}5$, a heuristic alignment-weighted scoring step
that adds a small consistent gain over that baseline (Exp.~3), and a positioning
run against eggNOG-mapper in which the embedding retrieval leads in most NK and
LK cells (Exp.~7). Those early figures predate the leakage fix and the cafaeval
protocol, so they are kept in the appendix for provenance only and carry no live
\texttt{/benchmark} or \texttt{/evaluation} source; a controlled comparison of
the baseline, the heuristic, and eggNOG-mapper on a shared leakage-clean
SELECT-window split is future work, covered by the score-ablation track. The
multi-PLM $\times$ $K$ grid (FARM-EXP.13) provides the empirical evidence base
for assessing whether the champion result holds across alternative PLM
representations.

\textbf{RQ2}: \emph{Do pairwise alignment statistics and taxonomic distance
provide complementary signals that improve the ranking of candidate annotations?}

Partly, and the honest answer is more nuanced than a single dominant family.
The development trace (\cref{app:dev-trace}) showed, before the leakage fix,
that layering pairwise alignment statistics on top of the embedding-only
baseline yields a small consistent gain (Exp.~3), directional evidence that
alignment carries signal partially orthogonal to the learned representation.
The leakage-free full-feature
importance analysis (\cref{tab:feature-importance}) then locates the dominant
ranking signal elsewhere: vote-aggregation features
(\texttt{neighbor\_vote\_fraction}, \texttt{go\_term\_frequency}) lead in every
category, GO-embedding similarity (\texttt{anc2vec\_neighbor\_cos},
\texttt{anc2vec\_query\_known\_maxcos}) contributes prominently as a
complementary ontology-aware term, and raw protein-embedding distance falls
outside the top-10, used by the model as a candidate-generation filter rather
than a final ranking signal. Classical alignment and taxonomy features rank
low individually once embeddings and neighbourhood statistics are present, so
their marginal contribution is real but modest. The answer to RQ2 is therefore
qualified: alignment and taxonomy add a small, genuine complementary signal,
while the decisive ranking information comes from the statistical structure of
the retrieved neighbourhood.

A methodological caveat governs how this question must be read, and it is
itself one of the rigor findings of the work. An earlier 56-feature
leave-one-family-out ablation (Exp.~8, \cref{tab:exp8-ablation}) appeared to
show the \texttt{anc2vec\_query} family dominating by a wide margin. That
magnitude does not reflect biological signal. It was traced to a
replication-by-category construction artefact in the CAFA-style training-pair
export: each (query, candidate) pair was replicated once per category bucket
(NK, LK, PK) and relabelled per bucket, so \texttt{anc2vec\_query\_known\_count}
effectively encoded the provenance of a row (a genuine positive versus a
synthetic negative injected from another bucket) rather than any ontology-derived
property. The model learned this provenance shortcut, and the inflated ablation
delta is its footprint, not evidence of feature value. The artefact was
diagnosed and patched in PROTEA (commit \texttt{223299c}); datasets built after
the fix, including the \texttt{bench-v1-K5-filtered} set behind
\cref{tab:feature-importance}, do not exhibit it. No primary claim in this
chapter rests on the pre-fix ablation magnitude. A clean re-derivation of which
features carry genuine signal, scored on full ground truth with
information-accretion weighting and free of the replication artefact, is the
object of the score-ablation study reported in the experimentation track as
future work.

\textbf{RQ3}: \emph{How can a reproducible, scalable annotation pipeline be
designed so that researchers can re-run analyses as ontologies and databases are
updated?}

PROTEA's versioned data model ensures reproducibility by design. Every
prediction is linked to an immutable \texttt{EmbeddingConfig} UUID, a specific
\texttt{AnnotationSet} (versioned by GOA release number), and a specific
\texttt{OntologySnapshot} (versioned by OBO release date). The temporal holdout
training pipeline is itself parameterised by version numbers: re-running the
same \texttt{train\_reranker\_auto} payload against the same database state
produces identical models and evaluation results. When new GOA releases become
available, the pipeline can be extended by appending the new version to the
training list and shifting the test window forward.

\subsection{The Value of Feature Engineering in the Age of Embeddings}

A recurring theme in this evaluation is the interplay between protein language
model embeddings and classical sequence analysis features. The embedding-only
baseline captures broad functional similarity through the learned representation
space, but its ranking is limited to a single signal: cosine distance.

The LightGBM re-ranker, when given access to the full 56-feature set, learns
to combine multiple complementary signals:
\begin{itemize}
  \item \textbf{Embedding distance} provides the initial KNN retrieval signal.
  \item \textbf{Alignment identity and scores} capture sequence-level conservation
    that the protein language model embedding may compress away.
  \item \textbf{Aggregation features} (vote fraction, GO term frequency) encode
    the statistical confidence of each candidate annotation.
  \item \textbf{GO-embedding similarity} (\texttt{anc2vec}) provides an
    ontology-aware signal complementary to protein-sequence embeddings.
  \item \textbf{Taxonomic distance} provides phylogenetic context.
  \item \textbf{Lineage features} (T-RES.1b; \texttt{lineage\_is\_ancestor\_of\_known},
    \texttt{lineage\_is\_descendant\_of\_known}, etc.) encode whether the candidate
    GO term is an ancestor or descendant of any term already known for the query
    protein.
\end{itemize}

The leakage-free full-feature importance analysis
(\cref{tab:feature-importance}) shows that vote-aggregation signals
(\texttt{neighbor\_vote\_fraction}, \texttt{go\_term\_frequency}) dominate
across all three categories, with GO-embedding similarity (\texttt{anc2vec})
contributing as a complementary ontology-aware term and raw protein-embedding
distance ranking outside the top-10. The much larger \texttt{anc2vec\_query}
ablation magnitude in the earlier 56-feature study
(\cref{tab:exp8-ablation}) is, as set out in the RQ2 answer above, a
replication-by-category dataset-construction artefact rather than evidence of
feature value; it is retained here as a methodological finding, patched in
PROTEA, and no conclusion in this section rests on it. With that caveat, the
leakage-clean evidence still supports the broader claim that feature
engineering remains valuable even when powerful embeddings are available: the
statistical structure of the retrieved neighbourhood and ontology-derived
similarity carry information not present in the raw protein-sequence embedding
alone. A clean per-feature re-derivation under full ground-truth
information-accretion scoring is provided by the score-ablation study reported
as future work.

\subsection{Limitations}
\label{subsec:limitations}

\begin{itemize}
  \item The temporal holdout strategy reduces contamination but does not
    eliminate it entirely: some test proteins share high sequence identity with
    reference proteins, making their annotations trivially predictable.
  \item ESMC~300M was chosen for computational tractability; larger variants
    would likely improve embedding quality at higher GPU cost.
  \item The evaluation considers only annotation transfer from Swiss-Prot
    experimental annotations; \gls{iea} annotations are excluded because they are
    themselves computationally derived.
  \item The IA file used (\texttt{IA\_cafa6.tsv}) was computed for the CAFA6
    challenge and may not perfectly reflect the information content of the
    specific GOA snapshots used here.
  \item The comparison with external tools is limited to eggNOG-mapper;
    additional tools (BLAST, InterProScan) would provide a more complete
    picture of the competitive landscape.
  \item The LK-BPO cell consistently favours the heuristic over the re-ranker,
    suggesting that the learned model may underperform for specific
    category-aspect combinations with limited positive examples.
  \item The Exp.~8 hyperparameter grid reveals a 0.0283 spread on NK-BPO
    (0.4386 to 0.4669), exceeding the informal robustness threshold of 0.02,
    indicating that optimal hyperparameters are cell-dependent and a joint
    per-cell tuning pass is warranted.
  \item The multi-PLM $\times$ $K$ generalisation grid (Exp.~9,
    \cref{tab:exp9-multiplm-grid}) is to be reported from the leakage-clean
    SELECT-window benchmark and is packaged via the F-DATA-PACK loop. The
    champion result (0.7291 $\pm$ 0.0028) is established on ProstT5 and does
    not depend on the grid; the grid provides breadth context for the
    generalisability claim.
\end{itemize}
